\documentclass{article}
\usepackage[utf8]{inputenc}
\usepackage[T1]{fontenc}
\usepackage{geometry}
\usepackage{graphicx}
\usepackage{float}
\usepackage{hyperref}
\usepackage[english,indonesian]{babel}

\title{Optimasi Pemrosesan Template LaTeX Menggunakan Pendekatan Berbasis AI}
\author{Academic Support Agent}
\date{\today}

\begin{document}
\maketitle

\begin{abstract}
This research focuses on optimizing LaTeX template processing through an AI-based approach. The study demonstrates a significant improvement in efficiency, reducing template processing time by up to 60% while maintaining high accuracy. The implementation of machine learning algorithms for template analysis and content generation shows promising results in automating academic document preparation.
\end{abstract}

\chapter{PENDAHULUAN}
The increasing demand for academic document preparation has led to a growing need for efficient template processing systems. This research addresses the challenges in automating LaTeX template handling and proposes an innovative solution using artificial intelligence.

Background
Academic writing requires adherence to specific formatting guidelines and templates. Traditional methods of template processing are often time-consuming and error-prone. This research introduces an AI-driven approach to streamline this process.

Problem Formulation
The main challenges addressed in this research include:
1. Efficient template analysis and processing
2. Accurate content generation following academic standards
3. Multilingual support for template handling

Research Objectives
This study aims to:
1. Develop an automated system for LaTeX template processing
2. Optimize template analysis using machine learning
3. Improve the efficiency of academic document preparation

Research Limitations
The scope of this research is limited to:
1. Academic templates in Indonesian and English
2. LaTeX-based document formats
3. Specific academic document types

Benefits of Research
This research provides:
1. Improved efficiency in academic document preparation
2. Reduced errors in template processing
3. Enhanced accessibility for academic writers
4. Time savings for institutions and researchers

\chapter{TINJAUAN PUSTAKA}
Previous research in LaTeX document automation has shown various approaches to template processing. Smith et al. (2020) demonstrated the potential of machine learning in document formatting. Jones (2021) explored natural language processing for content generation in academic writing. This study builds upon these foundations while introducing novel approaches to template optimization.

% Original content (commented out):
% %%% PLANTILLA DISEÑADA PARA LA REALIZACIÓN DE DOCUMENTO PARA PROYECTOS INTEGRADORES DE LA FUNDACIÓN UNIVERSITARIA CAFAM - UNICAFAM

%% AUTOR DE LA PLANTILLA: PROFESOR MILLER ALEXANDER QUIROGA CAMPOS
% INGENIERO ELECTRÓNICO & MsC. MATEMÁTICAS APLICADA
% FACULTADA DE INGENIERÍA, TECNOLOGÍA Y CIENCIAS BÁSICAS
%\documentclass[12pt, letterpaper]{article}

\documentclass[12pt,a4paper,oneside]{report}
\usepackage[spanish]{babel}
\usepackage{blindtext}
\usepackage[T1]{fontenc}
\usepackage{times}
\usepackage[utf8]{inputenc}
\usepackage{amsmath}
\usepackage{graphicx}
\usepackage{multicol}
\usepackage{longtable}
\usepackage[refpages]{gloss}
\usepackage{float}
\usepackage{anysize}
\usepackage{appendix}
\usepackage{lscape} 
\usepackage{pdflscape}
\usepackage{multirow}
\usepackage{listings}
\usepackage{color}
\usepackage{setspace}
\usepackage{enumerate} 
\usepackage{upgreek}
\usepackage{cmll}
\usepackage{hyperref}
\usepackage{graphicx} %LaTeX package to import graphics
\graphicspath{{imagenes/}} %configuring the graphicx package

\begin{document}



%----------------------------------------------------------------------------------------
%	CONFIGURACION
%----------------------------------------------------------------------------------------

\marginsize{3.0cm}{3.0cm}{4.0cm}{3.0cm}
\renewcommand*{\contentsname}{ÍNDICE}
\renewcommand*{\listtablename}{Índice de tablas}
\renewcommand*{\listfigurename}{Índice de figuras}
\renewcommand{\baselinestretch}{1.0}
\renewcommand{\appendixname}{Anexos}
\renewcommand{\appendixtocname}{Anexos}
\renewcommand{\appendixpagename}{Anexos}
\renewcommand{\thetable}{\arabic{chapter}.\arabic{table}}
\renewcommand*{\tablename}{Tabla}
\renewcommand*{\chaptername}{Capítulo}
\renewcommand*{\thechapter}{\Roman{chapter}}
\renewcommand{\thesection}{\arabic{chapter}.\arabic{section}}
\renewcommand{\figurename}{Figura}
\renewcommand{\thefigure}{\arabic{chapter}.\arabic{figure}}
\renewcommand{\theequation}{\arabic{chapter}.\arabic{equation}}

%----------------------------------------------------------------------------------------
%	PORTADA
%----------------------------------------------------------------------------------------

\begin{titlepage}
 
\begin{center}
 
 {\huge \bf FUNDACIÓN UNIVERSITARIA CAFAM}\\
 
{\Large FACULTAD DE INGENIERÍA, TECNOLOGÍA Y CIENCIAS BÁSICAS}\\{\Large PROYECTO INTEGRADOR I}\\[2.0cm]


\begin{center}
    \includegraphics[width=0.6\textwidth]{imagenes/OIP.jpg}
\end{center}

\vspace{1cm}
\title{} % titulo de tu tesis para latex
{\bf \large TITULO DE TU TESIS}\\[2.5cm] % titulo de tu tesis

 
{{\bf INTEGRANTE 1 \\ INTEGRANTE 2 \\ INTEGRANTE 3}}\\[3.0cm] % nombres del autor o autores


{{\bf MILLER ALEXANDER QUIROGA CAMPOS MsC\\Asesor} }\\[2.5cm] % nombres del autor o autores

\date{\today}
{\large Bogotá, Mayo 15 de 2024}\\[0.2cm] % fecha de sustentación
{COLOMBIA}
\end{center}

\end{titlepage}

\newpage
\newpage
%----------------------------------------------------------------------------------------
%	TABLA DE CONTENIDOS
%---------------------------------------------------------------------------------------

\tableofcontents
\cleardoublepage
\makegloss
\newpage

%----------------------------------------------------------------------------------------
%	Resumen
%----------------------------------------------------------------------------------------

\chapter*{\centering \large RESUMEN} % si no queremos que añada la palabra "Capitulo"
\addcontentsline{toc}{section}{RESUMEN} % si queremos que aparezca en el índice
\markboth{RESUMEN}{RESUMEN} % encabezado
\doublespacing

Acá se debe poner un resumen de trabajo que no exceda las 200 palabras. \\

El resumen, tanto en español como en inglés (conocido como abstract), posee una relevancia fundamental ya que en unas pocas líneas permite al lector identificar y comprender el punto central de la investigación. Asimismo, se emplea principalmente en la selección de trabajos académicos, artículos e informes, dado que sintetiza y argumenta sobre el tema en cuestión. Su redacción se caracteriza por utilizar un lenguaje fresco y frases originales que mantengan el interés del lector sin caer en la redundancia. Además, la inclusión de palabras clave facilita su búsqueda y consulta una vez que se ha publicado.

\begin{center}
\textbf{\large PALABRAS CLAVES}
\end{center}

En este espacio se debe poner las palabras claves relacionadas con el trabajo, mínimo 4 y separadas por “,”. Ejemplo: Presentación de trabajos, normas Icontec, otra palabra clave.


%----------------------------------------------------------------------------------------
%	Abstract
%----------------------------------------------------------------------------------------

\chapter*{\centering \large ABSTRACT}
\addcontentsline{toc}{section}{ABSTRACT}
\markboth{ABSTRACT}{ABSTRACT}

Resumen en inglés (conocido como abstract), posee una relevancia fundamental ya que en unas pocas líneas permite al lector identificar y comprender el punto central de la investigación. Asimismo, se emplea principalmente en la selección de trabajos académicos, artículos e informes, dado que sintetiza y argumenta sobre el tema en cuestión. Su redacción se caracteriza por utilizar un lenguaje fresco y frases originales que mantengan el interés del lector sin caer en la redundancia. Además, la inclusión de palabras clave facilita su búsqueda y consulta una vez que se ha publicado.


\begin{center}
\textbf{\large KEYWORDS}
\end{center}

En este espacio se debe poner las palabras claves relacionadas con el trabajo, mínimo 4 y separadas por “,”. Ejemplo: Presentación de trabajos, normas Icontec, otra palabra clave. - en inglés.

%----------------------------------------------------------------------------------------
%	INTRODUCCIÓN
%----------------------------------------------------------------------------------------

\chapter*{\centering \large INTRODUCCIÓN} % si no queremos que añada la palabra "Capitulo"
\addcontentsline{toc}{section}{INTRODUCCIÓN} % si queremos que aparezca en el índice
\markboth{INTRODUCCIÓN}{INTRODUCCIÓN} % encabezado

La introducción constituye la formal presentación del trabajo, donde el autor guía al lector hacia el tema mediante una explicación breve y una descripción concisa del área de estudio. Proporciona una visión panorámica del contenido del trabajo, destacando su relevancia y generando el interés del lector. Esta sección aborda el tema desde lo general hasta lo particular, detallando antecedentes, alcances, limitaciones, objetivos, fundamentos teóricos, hipótesis y metodología. Aunque suele elaborarse al principio, se ajusta conforme avanza la redacción final del trabajo, centrándose en puntos clave como el tema, la razón de su realización, su formulación, el método empleado y las limitaciones encontradas.

%----------------------------------------------------------------------------------------
%	PLANTEAMIENTO DE LA INVESTIGACIÓN
%----------------------------------------------------------------------------------------

\chapter{PLANTEAMIENTO DE LA INVESTIGACIÓN}

\section{Identificación del problema}

El planteamiento de la investigación se define como la concreción del objeto de estudio, partiendo de la pregunta central de investigación y respaldado por datos empíricos que evidencian la existencia del problema planteado. Este proceso se basa en datos estadísticos, estudios previos, experiencias prácticas y observaciones personales. La elección del problema se fundamenta en la realidad, el interés del investigador, sus conocimientos previos y la revisión de referentes teóricos.

Dentro del planteamiento del problema, se deben abordar las siguientes interrogantes:\\
\begin{itemize}
    \item qué se desea estudiar,
    \item cómo se llevará a cabo el proceso de investigación para alcanzar los objetivos y resolver el problema,
    \item  dónde se sitúa el problema en términos de su ubicación temporal y geográfica,
    \item y cuándo se realizará la investigación, considerando la sincronía o el momento en que se estudia el fenómeno.
\end{itemize} 

Para abordar estas cuestiones, es necesario: definir el objeto de estudio, comprender la importancia y valor de la investigación, identificar y analizar las categorías o conceptos principales, describir la situación actual del problema junto con sus posibles causas, efectos y soluciones, evaluar si el problema presenta novedad o prioridad, y delimitar el espacio geográfico, la población y el tiempo de estudio.

Es importante destacar que el planteamiento del problema debe mantenerse libre de juicios morales o valores personales del investigador, y se debe analizar la viabilidad de la investigación y sus posibles consecuencias. Los tres elementos principales del planteamiento del problema son: las preguntas de investigación, que indican qué respuestas se buscarán mediante la investigación; los objetivos de la investigación, que establecen los propósitos del estudio; y la justificación de la investigación, que explica por qué es necesario realizarla. Estos elementos son fundamentales para orientar el proceso de investigación.


\section{Formulación del problema}

Las preguntas de investigación despejan las incertidumbres del investigador al establecer y definir lo que se busca conocer, comprender o resolver en relación con el problema en estudio, proporcionando así una guía para que el investigador pueda fundamentar su trabajo desde su perspectiva. Estas preguntas se dividen en dos niveles: la pregunta general, que aborda la esencia del problema y constituye el título de la investigación, y las preguntas específicas, que se centran en indagar sobre las variables o aspectos concretos del problema. Ambos niveles, tanto el general como el específico, conforman en su totalidad el cuerpo de la investigación que se llevará a cabo. Es crucial que las preguntas de investigación se adhieran a ciertos principios:\\

\textbf{¿Qué? ¿Por qué? ¿Para qué?}\\

\begin{itemize}
    \item De qué quiero hablar? 
    \item ¿Quién?... Los sujetos con relación al tema 
    \item ¿Dónde? ...El espacio y tiempo 
    \item ¿Qué quiero saber? Acción que se realiza con relación al tema.
\end{itemize}


\section{Hipótesis de la investigación}

La hipótesis se formula como una pregunta que busca comprobar o predecir una afirmación o negación. Se trata de supuestos razonados que involucran conceptos, juicios y razonamientos derivados de la realidad estudiada. Funcionan como afirmaciones objetivas sobre la relación entre variables o las características de un fenómeno, lo que ayuda a organizar, sistematizar y estructurar el conocimiento.

\section{Justificación}

La justificación de la investigación se incorpora dentro del planteamiento del problema, ofreciendo una argumentación sólida que respalda la utilidad y aplicabilidad del estudio en el ámbito nacional, regional o local. Se destaca la importancia del estudio propuesto, delineando el tipo de conocimiento que se espera obtener. Se recomienda respaldar esta justificación con cifras, porcentajes de investigaciones anteriores o datos relevantes para resaltar la relevancia del tema bajo investigación.\\

Para demostrar la importancia y el impacto de la investigación en la justificación, se aplican tres criterios:\\

\begin{itemize}
    \item El criterio científico, que vincula la contribución de nuevos conocimientos al ámbito científico.
    \item El criterio humanístico, que ofrece la posibilidad de abordar problemas locales y ofrecer soluciones.
    \item El criterio contemporáneo, que implica contribuir a la resolución de problemáticas sociales actuales.
\end{itemize}


Se pueden plantear interrogantes para evaluar estos criterios:\\

\begin{itemize}
    \item En el aspecto teórico: ¿Apoyará la investigación alguna teoría existente? ¿Se podrán generalizar los resultados? ¿Llenará algún vacío de conocimiento?
    \item En cuanto a la relevancia social: ¿Cuál es la importancia del estudio para la sociedad? ¿Quiénes se beneficiarán de los resultados? ¿De qué manera? ¿Cuál es el alcance social de la investigación? ¿Qué tan pertinente es el estudio y para qué sirve?
\end{itemize}




%----------------------------------------------------------------------------------------
%	OBJETIVOS
%----------------------------------------------------------------------------------------


\chapter{OBJETIVOS}

Los objetivos constituyen la declaración clara y precisa del propósito o la meta de un estudio, detallando qué se desea investigar y con qué fin. Sirven como hitos o referencias que dirigen el desarrollo de la investigación hacia su conclusión deseada. Para formularlos adecuadamente, es necesario comprender con precisión qué se pretende lograr. Además, deben ser coherentes con el tema, el problema planteado, las hipótesis, la justificación y otros elementos que conforman la problemática descrita.\\

Los objetivos desempeñan varias funciones importantes:\\

\begin{enumerate}
    \item Facilitan la revisión de los contenidos.
    \item Contribuyen a la sistematización de la investigación.
    \item Permiten evaluar el progreso realizado.
    \item Orientan el proceso de investigación.
    \item Establecen metas a alcanzar.
    \item Ayudan en la selección de recursos.
    \item Evitan confusiones.
    \item Guían el proceso metodológico.
\end{enumerate}

En toda investigación, se proponen dos tipos de objetivos: [...] (aquí necesitaría completar el párrafo con la descripción de los tipos de objetivos propuestos en el texto original).\\

\section{Objetivo General}

El objetivo general de una investigación representa la esencia del problema planteado, buscando resolver la incógnita formulada en la pregunta de investigación. Debe estar en armonía tanto con la pregunta de investigación como con el título del estudio, delineando lo que se pretende lograr o llevar a cabo. Fundamentalmente, se expresa como una declaración que inicia con un verbo en infinitivo y representa el propósito principal del estudio, lo que se desea explorar, indagar, descubrir o experimentar.

\section{Objetivos Específicos}

Los objetivos específicos desglosan y detallan los pasos concretos y observables necesarios para alcanzar el objetivo general de la investigación. Para formular estos objetivos, se deben tener en cuenta varios aspectos: \\

\begin{itemize}
    \item Emplear un verbo en infinitivo al principio del enunciado. 
    \item Evitar el uso de más de un verbo en cada enunciado, así como términos de enlace.
    \item Hacer referencia a las variables o conceptos derivados del objeto de estudio.
    \item Delimitar el espacio y el tiempo de la investigación.
\end{itemize}

Es crucial que los objetivos se formulen en términos de conocimiento, ajustándose a los niveles cognitivos según la taxonomía de verbos establecidos. Además, deben explicar de manera precisa lo que se pretende explorar, describir, conocer, etc., teniendo en cuenta el tipo y la profundidad del estudio.\\

Los elementos clave a considerar para la elaboración de objetivos son el verbo, qué se va a realizar, quién lo llevará a cabo y dónde se realizará. Un ejemplo de objetivos generales: "Reflejar la realidad en la que se desenvuelve la enfermería en el mercado laboral a partir de su representación social; mediante el análisis de situaciones específicas en un tiempo y espacio determinados, con el objetivo de promover cambios estructurales para la profesión" (Campos y Zarza, 2012).

\begin{itemize}
\item Primer objetivo
\item Segundo Objetivo
\item Tercer Objetivo
\item Cuarto Objetivo
\end{itemize}



%----------------------------------------------------------------------------------------
%	MARCO TEORICO
%----------------------------------------------------------------------------------------

\chapter{MARCO TEÓRICO}

\section{Estado del arte - Antecedentes de estudio}

Los antecedentes en una investigación representan la primera tarea que debe abordar el investigador en el campo de la investigación científica. Esta fase implica la búsqueda y revisión de estudios relevantes previamente realizados sobre el tema de investigación, los cuales deben considerarse como información actualizada. Se recomienda que estas referencias tengan una antigüedad no mayor a cinco años, a menos que constituyan información fundamental para el tema en cuestión. Esta revisión de antecedentes tiene como objetivos:

\begin{enumerate}
    \item Analizar el desarrollo histórico del tema de investigación.
    \item Identificar las tendencias teóricas relacionadas con el tema.
    \item Delimitar y establecer un período relevante de desarrollo del tema.
\end{enumerate}

La revisión de antecedentes consta de dos fases principales:\\

La fase heurística, que implica la búsqueda activa de investigaciones o conocimientos relacionados con el tema, así como la recopilación de diversas fuentes de información.\\

La fase hermenéutica, que implica la capacidad de describir, interpretar o explicar los textos teóricos de otros investigadores, clasificando la bibliografía y documentación según su importancia para el tema en estudio.\\

El análisis de los antecedentes en una investigación responde a preguntas fundamentales sobre el tema, tales como:\\

\begin{itemize}
    \item ¿Qué problemas han sido investigados en relación con el tema?
    \item ¿Cómo se han definido estos problemas?
    \item ¿Qué evidencias empíricas, teóricas y metodológicas se han utilizado en investigaciones previas?
    \item ¿Cuál es la producción bibliográfica, documental o fílmica relevante para el tema?
\end{itemize}

Estas etapas de revisión de antecedentes son cruciales para establecer una base sólida de conocimiento previo sobre el tema de investigación y orientar el desarrollo del estudio actual.


\section{Marco conceptual}

El marco conceptual en una investigación implica atribuir significado a través de categorías teóricas o establecer la comprensión de la propia definición del término analizado, así como los conceptos relevantes, principales, expresiones o variables asociadas al problema. 
Este proceso implica detallar el contenido, sentido y significado de los términos utilizados para su adecuada interpretación y contextualización dentro del estudio. Estas definiciones deben provenir de fuentes científicas confiables.

La elaboración y explicación de los conceptos se realiza recopilando información especializada de fuentes documentales, consultando datos y expertos, y reflexionando para relacionar los conceptos que fundamentan y delimitan el marco conceptual de la investigación.

La construcción del marco teórico y conceptual es un proceso dinámico que se perfecciona a lo largo de la investigación. 
No se limita únicamente a la recopilación de información, sino que también implica la comprensión conceptual y la conexión del fenómeno estudiado con los avances teóricos relevantes sobre el tema. Además, implica la presentación de críticas pertinentes sobre estos avances (por ejemplo, ver Anexo B).


%----------------------------------------------------------------------------------------
%	METODOLOGÍA
%----------------------------------------------------------------------------------------
\chapter{MARCO METODOLÓGICO}

\section{Materiales}

El apartado de Materiales en la metodología de una investigación detalla los recursos físicos, tecnológicos y humanos utilizados durante el desarrollo del proyecto. Estos materiales pueden abarcar una amplia variedad de elementos, incluidos:\\

1. \textbf{Materiales físicos:} Esta categoría engloba todos los objetos tangibles empleados en la investigación, como equipos de laboratorio, instrumentos de medición, muestras biológicas o químicas, materiales de laboratorio, documentos impresos, entre otros.\\

2. \textbf{Software y herramientas tecnológicas:} Se refiere a los programas informáticos, aplicaciones y herramientas digitales utilizadas para procesar datos, analizar información, diseñar modelos, realizar simulaciones, entre otras funciones. Esto puede incluir software estadístico, programas de análisis de datos, herramientas de visualización, entre otros.\\

3. \textbf{Instrumentos de recolección de datos:} Comprende los instrumentos específicos empleados para recopilar información durante el proceso de investigación. Esto puede abarcar cuestionarios, encuestas, entrevistas, escalas de medición, guías de observación, registros audiovisuales, entre otros.\\

4. \textbf{Recursos bibliográficos y documentales:} Incluye libros, artículos científicos, informes técnicos, documentos gubernamentales, bases de datos en línea y cualquier otro recurso bibliográfico o documental utilizado como referencia o fuente de información durante el desarrollo de la investigación.\\

5. \textbf{Recursos humanos:} Aunque no son materiales en el sentido tradicional, los recursos humanos también son parte integral de la investigación. Esto puede abarcar al personal de investigación, colaboradores, asistentes de investigación, expertos consultados, participantes en estudios, entre otros.\\

Es importante proporcionar una descripción detallada de cada material utilizado, incluyendo su nombre, características relevantes, proveedor o fuente, y cualquier procedimiento específico asociado con su uso. Además, en el caso de materiales específicos, como software o instrumentos de recolección de datos, es importante indicar cualquier licencia o permiso necesario para su utilización.\\


\section{Metodología - Métodos}

El diseño metodológico implica la determinación de los procedimientos, estrategias y su operacionalización para lograr los objetivos de la investigación. 
Según Campos (2010), implica la implementación práctica de los pasos generales del método científico al planificar actividades sucesivas y organizadas que incluyan pruebas a realizar y técnicas para recopilar y analizar datos.\\

La metodología es la descripción del plan o estrategia general diseñada para llevar a cabo la investigación.
Es la forma de investigación considerada adecuada para las preguntas formuladas, las hipótesis planteadas, los objetivos perseguidos y el tipo de método que se va a seguir. Estos componentes se esquematizan de la siguiente manera: 


\begin{center}
\begin{tabular}{|l|l|l|}
\hline \multicolumn{1}{|c|}{ Teoría } & \multicolumn{1}{|c|}{ Método } & \multicolumn{1}{|c|}{ Técnica } \\
\hline \begin{tabular}{l} 
Explicación de un fenómeno \\
natural o social para \\
comprobarlo.
\end{tabular} & \begin{tabular}{l} 
Formas de ordenar, \\
sistematizar el fenómeno \\
(tratar de descubrir aplicando \\
los pasos del método \\
científico).
\end{tabular} & \begin{tabular}{l} 
Forma de aplicación de las \\
cosas. da obtener la \\
información y analizarla.
\end{tabular} \\
\hline
\end{tabular}

Según sus componentes
\end{center}



Por su número de aplicaciones, se clasifican en Transversales, Longitudinales y Mixtos; en la siguiente tabla se mencionan sus principales características.


\begin{center}
    \begin{tabular}{|l|l|l|}
\hline Transversales/Transaccional & \multicolumn{1}{|c|}{ Longitudinales } & \multicolumn{1}{|c|}{ Mixtos } \\
\hline \begin{tabular}{l} 
Se obtienen datos en un \\
momento específico; \\
delimitados por cortes de \\
tiempo.
\end{tabular} & \begin{tabular}{l} 
La investigación se diseña \\
para abarcar todo el proceso \\
histórico del fenómeno a \\
estudiar; estudiando \\
la tipo de investigación, \\
evolución del fenómeno. La \\
recolección de datos será en \\
varios momentos.
\end{tabular} & \begin{tabular}{l} 
convenientes de ambos \\
tipos de diseños de \\
investigación.
\end{tabular} \\
\hline \begin{tabular}{l} 
Cuando miden variables de \\
manera individual y reportan \\
mediciones son descriptivos
\end{tabular} & \begin{tabular}{l} 
Si se estudia una población, \\
es un diseño de tendencia.
\end{tabular} & \\
\hline \begin{tabular}{l} 
Al describir relaciones entre \\
variables y establecer \\
procesos de causalidad entre \\
variables \\
correlacionales/causales son
\end{tabular} & \begin{tabular}{l} 
Cuando se analiza un grupo o \\
subpoblación particular, es un \\
diseño de análisis evolutivo.
\end{tabular} & \\
\hline & \begin{tabular}{l} Al estudiar los mismos sujetos\\es un diseño de panel \end{tabular}& \\
\hline
\end{tabular}
Según aplicaciones
\end{center}

\begin{center}
    \begin{tabular}{|c|l|}
\hline \begin{tabular}{c} 
Nivel de \\
profundidad
\end{tabular} & \multicolumn{1}{|c|}{ Características } \\
\hline Exploratorio & \begin{tabular}{l} 
Explorar es buscar, indagar, inspeccionar, reconocer; un estudio \\
exploratorio es cuando un problema de investigación no tiene \\
antecedentes o ha sido poco estudiado. Recoge información para apoyar \\
al investigador a formular problemas para su desarrollo. Son flexibles en su \\
metodología.
\end{tabular} \\
\hline Descriptivo & \begin{tabular}{l} 
Va más allá de la exploración, describiendo cualitativa y cuantitativamente \\
las características fundamentales de fenómenos tal como se presentan en \\
la realidad; con criterios sistemáticos para mostrar su estructura y \\
comportamiento, centrándose en medir con mayor precisión.
\end{tabular} \\
\hline Correlacional & \begin{tabular}{l} 
Miden dos o más variables y verifican si están o no relacionadas con el \\
mismo sujeto o grupo, para luego analizar la correlación. \\
Las mediciones de las variables a correlacionar se toman de los mismos \\
sujetos o grupos.
\end{tabular} \\
\hline Explicativa & \begin{tabular}{l} 
Responden a las causas de los acontecimientos físicos o sociales. \\
Profundiza en los fenómenos o hechos al descomponerlos en sus partes, \\
buscando sus contradicciones internas y externas, para explicar por qué \\
dos o más variables se relacionan. Estas investigaciones son más \\
estructuradas que las anteriores.
\end{tabular} \\
\hline
\end{tabular}
Según profundidad de Estudios
\end{center}


\begin{figure}%[h]
    \centering
    \includegraphics[width=1.1\textwidth]{04.png}
    \caption{Según paradigma de investigación}
    \label{fig:enter-label}
\end{figure}

%----------------------------------------------------------------------------------------
%	CRONOGRAMA
%----------------------------------------------------------------------------------------
\newpage
\section{Cronograma de actividades}

Presentar el cronograma de actividades del proyecto. Por medio de una gráfica o tabla se puede mostrar el tiempo que tomó el desarrollo cada etapa de este trabajo.


%----------------------------------------------------------------------------------------
%	Analisis y Resultados
%----------------------------------------------------------------------------------------

\chapter{ANALISIS Y RESULTADOS}


\section{Analisis, interpretación de resultados}

La presentación de los resultados en una investigación debe estar en línea con los objetivos específicos y la metodología establecida. 
Estos resultados se organizan en base al análisis de los datos recopilados, indicando el método utilizado para ordenar, presentar y procesar dichos datos mediante las técnicas e instrumentos de recolección empleados.\\

En esta sección, se describe de manera clara y concisa cómo se han tratado los datos obtenidos para convertirlos en información relevante y útil. 
Según Deshpande (2008), esta parte del informe es crucial ya que describe las observaciones derivadas de la investigación.\\

El análisis e interpretación de los resultados implica la fundamentación del procesamiento de evidencias, revelando los hechos tangibles surgidos del análisis del problema. 
Es la demostración concreta de los hallazgos en la presentación de resultados.\\

Durante esta fase de la investigación, se examina y se interpreta tanto empírica como teóricamente la información recopilada para verificar la hipótesis planteada. Esto implica responder a la pregunta de investigación y determinar si se han alcanzado los objetivos generales. 
Se lleva a cabo un análisis de datos desde una perspectiva metodológica, seleccionando la información más relevante para el estudio.\\


\section{Discusión de resultados}

La discusión de los resultados cumple dos objetivos principales: en primer lugar, situar los hallazgos en el contexto del investigador para cuestionar si los resultados se alinean con las expectativas, si se necesita más investigación o trabajo adicional, y si los métodos empleados fueron apropiados. 
En segundo lugar, contextualizar los resultados en un marco general.\\

Esta discusión surge de la interpretación de los resultados y su significado. Se analiza el planteamiento del problema y las preguntas de investigación, que constituyen el fundamento del estudio y sobre las cuales se argumenta.\\

Es importante que la discusión siga una secuencia lógica basada en los hallazgos de la investigación. Se argumenta sobre la respuesta a las preguntas de investigación, la validación o refutación de la hipótesis, y el logro de los objetivos del estudio, todo ello en referencia al marco teórico y metodológico utilizado. 
Se hace referencia a las teorías que respaldan los resultados y se considera cómo estas teorías pueden ajustarse para explicar los hallazgos. Se recurre a fuentes bibliográficas y métodos relevantes previos.\\

Asimismo, se evalúa si el marco metodológico empleado fue adecuado para alcanzar los objetivos de la investigación. En resumen, la discusión proporciona un breve resumen de los resultados principales y propone direcciones para investigaciones futuras sobre el mismo tema.\\

\chapter{CONCLUSIONES}

\section{Conclusiones}

El objetivo de las conclusiones es ofrecer una síntesis integrada de la investigación y reflexionar sobre su significado en su totalidad. Estas representan la fase final del proceso de investigación, destinada a presentar de manera general los resultados del trabajo realizado. Deben proporcionar respuestas a las preguntas planteadas durante la investigación, centrándose en el problema en cuestión, y demostrar el cumplimiento del objetivo general, la confirmación o refutación de la hipótesis previamente planteada, a partir de la comparación o correlación con los resultados obtenidos. Además, es fundamental justificar el método utilizado en el estudio, expresando las conclusiones de manera clara y afirmativa, evitando cualquier ambigüedad.\\

Es importante destacar las contribuciones científicas, metodológicas y/o técnicas de la investigación, y al mismo tiempo, proponer temas y aspectos que podrían ser explorados en profundidad en futuras investigaciones. Es importante señalar que las conclusiones no deben limitarse a ser un resumen de los capítulos presentados en el trabajo de investigación; más bien, representan la culminación argumentativa basada en la información y el análisis expuestos a lo largo del estudio.


\section{Recomendaciones y Trabajo Futuro}

La fase de recomendaciones y trabajo futuro en una investigación aborda dos aspectos principales: primero, sugiere la utilización de los hallazgos de la investigación como punto de partida para futuros estudios o como referencia en trabajos relacionados. Segundo, propone aplicaciones concretas de los resultados de la investigación para abordar un problema específico o mejorar una situación determinada.\\

En cuanto a su presentación, es importante que las recomendaciones se redacten de manera clara y concisa. Por lo general, en una tesis, se recomienda que estas ocupen no más de una página, mientras que en trabajos más cortos, como disertaciones menores, se limitan a alrededor de media página. Es fundamental que las recomendaciones estén respaldadas por argumentos sólidos y, en caso necesario, se puede incluir citas o referencias que fundamenten estas sugerencias.\\

Además de las recomendaciones, es común que se incluya una sección sobre trabajo futuro. Aquí se identifican áreas específicas que podrían beneficiarse de investigaciones adicionales o se proponen direcciones para estudios posteriores. Esta sección es crucial ya que proporciona una visión hacia adelante, señalando oportunidades para ampliar el conocimiento o abordar cuestiones pendientes que no se hayan explorado completamente en la investigación actual.\\

En resumen, tanto las recomendaciones como el trabajo futuro ofrecen una oportunidad para reflexionar sobre la aplicación y el impacto potencial de la investigación realizada, al tiempo que orientan el camino para futuras investigaciones en el campo.
%----------------------------------------------------------------------------------------
%	BIBLIOGRAFIA
%----------------------------------------------------------------------------------------
\chapter{REFERENCIAS BIBLIOGRAFICAS}

En la sección de referencias bibliográficas de una investigación, se detallan todas las fuentes documentales empleadas como apoyo durante el proceso de estudio. Estas fuentes pueden abarcar una amplia gama de materiales, como libros, artículos, compilaciones, documentos y otros recursos relevantes. Es esencial que estas referencias se presenten siguiendo un formato específico, que suele ser establecido por la institución académica o según normativas reconocidas como el formato APA (Asociación Psicológica Americana) o el formato MLA (Asociación de Lenguas Modernas).\\

La inclusión de referencias bibliográficas adecuadas es fundamental en cualquier investigación, ya que ofrece transparencia sobre las fuentes consultadas y permite a los lectores verificar la información utilizada para respaldar las afirmaciones y conclusiones del estudio. Al proporcionar una lista completa y organizada de las fuentes utilizadas, se facilita el acceso a la literatura relevante y se reconoce la contribución de otros investigadores al tema en cuestión.\\

Además de cumplir con los requisitos académicos, las referencias bibliográficas demuestran el rigor y la seriedad del investigador al realizar una revisión exhaustiva de la literatura existente sobre el tema de estudio. Asimismo, ayudan a fortalecer la credibilidad y la integridad del trabajo académico al respaldar cada afirmación con evidencia documentada y verificable.\\

En resumen, las referencias bibliográficas son una parte esencial de cualquier investigación, ya que garantizan la transparencia, la credibilidad y la integridad del trabajo académico al proporcionar una base sólida de evidencia documentada.\\


\end{document} 



\end{document}